\documentclass[11pt,letterpaper]{article}

\usepackage{graphicx}           % Graphics package
\usepackage{amsmath} 
\usepackage{mathrsfs}
\usepackage[colorlinks=true]{hyperref} 
\usepackage{color}
\usepackage{amsfonts}
\usepackage[textheight=9in, textwidth=7.5in, letterpaper]{geometry}
\usepackage[makeroom]{cancel}
\usepackage{cite}
\usepackage{sectsty}
\usepackage{empheq}
\usepackage{appendix}
\usepackage{enumerate} 
\usepackage{simpler-wick}
\usepackage{amssymb}

\sectionfont{\fontsize{12}{12}\selectfont}
\subsectionfont{\fontsize{12}{12}\selectfont}

\newcommand*\widefbox[1]{\fbox{\hspace{2em}#1\hspace{2em}}}
\newcommand{\LP}{\left(}
\newcommand{\RP}{\right)}
\newcommand{\mc}[1]{\mathcal{#1}}
\newcommand{\R}{\mathbb{R}}
\newcommand{\C}{\mathbb{C}}
\newcommand{\Z}{\mathbb{Z}}
\newcommand{\D}{\partial}
\newcommand{\KET}[1]{\left| #1 \right\rangle }
\newcommand{\BRA}[1]{\left\langle #1 \right| }
\newcommand{\IP}[2]{\left\langle #1 \middle| #2 \right\rangle}
\newcommand{\PA}[2]{\frac{\partial #1}{\partial #2}}



%%%%%%%%%%%%%%%%---------------------MOST USEFUL COMMAND EVER---------------------%%%%%%%%%%%%%%%%%%%%%%
\newcommand{\fixme}[1]{{\bf {\color{red}[#1]}}}
\newcommand{\draftmode}{\usepackage[notref,notcite]{showkeys}}
\providecommand*\showkeyslabelformat[1]{\normalfont\sffamily\footnotesize#1}

%%%%%%%%%%%%%%%


\setcounter{tocdepth}{4}
\setcounter{secnumdepth}{4}

\begin{document}

\title{Gravitation $w_{1 + \infty}$ Symmetry}

\author{Mathew Calkins}

\date{\today}

 
\maketitle

\abstract{Extremely rough notes on gravitation $w_{1 + \infty}$ symmetry. Maybe these will turn into something coherent, maybe not!}

\tableofcontents




\section{$w_{1 + \infty}$ and the Celestial Sphere}
Scratch notes following \cite{Strominger:2021mtt}.


\subsection{Gravity}
Denote by $G^+_\Delta(z, \bar{z})$ the celestial dual to a positive-helicity\footnote{TODO: Comment on choice to restrict to positive-helicity sector.} bulk graviton. The authors of \cite{Guevara:2021abz} find by analyzing bulk collinear limits that two such operators have OPE \begin{equation}\label{gravitonOPE}
G_{\Delta_1}^+ (z_1, \bar{z}_1) G_{\Delta_2}^+ (z_2, \bar{z}_2) = - \frac{\kappa}{2} \frac{1}{z_{12}} \sum_{n =0 }^\infty B(\Delta_1 - 1 + n, \Delta_2 - 1) \frac{\bar{z}_{12}^{n+1}}{n!} \bar{\D}^n G^+_{\Delta_1 + \Delta_2} (z_2, \bar{z}_2) + \mc{O}\LP z_{12}^0 \RP.
\end{equation}
Here $B$ is the Beta function \begin{equation}
B(x,y) = \frac{\Gamma(x) \Gamma(y)}{\Gamma(x+y)}
\end{equation}
holomorphic poles\footnote{TODO: Comment more on this choice.}. These operators have conformal spin $s = +2$ and scaling weights $\Delta$, hence conformal weights \begin{equation}
\LP h, \bar{h} \RP = \LP \frac{\Delta + s}{2}, \frac{\Delta - s}{2} \RP = \LP \frac{\Delta}{2} + 1, \frac{\Delta}{2} - 1 \RP.
\end{equation}
These operators are, by definition, the result of Mellin-transformating a bulk scattering amplitude of gravitons, necessarily in an asymptotically flat spacetime, to produce a boundary correlation function. \begin{equation}
\langle G_\Delta^\pm(z, \bar{z}) \cdots \rangle \equiv \int_0^\infty \frac{d \omega}{\omega} \omega^\Delta \cdots \mc{A} \LP \omega \hat{p}(z, \bar{z}),  \cdots \RP.
\end{equation}
From the well-known property of Mellin transformations \begin{equation}
\int_0^\infty \frac{d \omega}{\omega} \omega^\Delta \LP \lim_{\omega \to 0} \omega f(w) \RP = \lim_{\Delta \to 1} (\Delta - 1) \int_0^\infty \frac{d \omega}{\omega} \omega^\Delta \omega f(w)
\end{equation}
we see that a bulk amplitude involving a soft graviton \begin{equation}
\lim_{\omega \to 0} \omega \mc{A} \LP \omega \hat{p}(z, \bar{z}),  \cdots \RP
\end{equation}
corresponds on the boundary to taking the limit \begin{equation}
\lim_{\Delta \to 1} (\Delta - 1) G_\Delta^+(z, \bar{z}).
\end{equation}
This nicely agrees with \eqref{gravitonOPE}, which we see has an order-one pole at $\Delta_i = 1$, rendering this limit finite. In fact we see it has a tower of order-one poles at $\Delta_i = 1, 0, -1, \ldots$. We therefore define more generally \begin{equation}
H^k(z, \bar{z}) \equiv \lim_{\varepsilon \to 0} \varepsilon G_{k + \varepsilon}^+(z, \bar{z}), \ k \in \Z.
\end{equation}
We will shortly find that these live in a shortened representation, so that only a subset of $k$ values give a nonzero result.\\
\\
\
These objects turn out to generate an interesting algebra. We first expand in the antiholomorphic coordinate \begin{equation}
H^k(z, \bar{z}) = \sum_n \frac{H_n^k(z)}{\bar{z}^{n + \frac{k-2}{2}}}
\end{equation}
where the shift in the powers of $\bar{z}$ is required by $H^k$ having a nonzero antiholomorphic conformal weight. Again, we will find that the relevant values of $n$ are limited. We now interpret the objects $H_n^k(z)$ as holomorphic currents on the boundary. 


\bibliography{bibfile}
\bibliographystyle{utphys}

\end{document}


